\subsection{INNER PRODUCTS: CASE OF COGEBRAS}

Let \(E = \bigoplus_{p \geq 0} E_p\) be a coassociative counital graded cogebra. Then we know (no. 2, Propositions 1 and 3) that the graded dual \(E^{*gr}\) has (with the convention on graduations made at the beginning of the paragraph) a graded algebra structure of type N over A, the product of two elements \(u, v\) of this algebra being defined by \(uv = m \circ (u \otimes v) \circ c\) , where \(c: E \to E \otimes_A E\) is the coproduct and \(m: A \otimes_A A \to A\) defines the multiplication. In other words, if, for \(x \in E\) , \(c(x) = \sum_i y_i \otimes z_i\) , we can write (canonically identifying \(A \otimes_A E\) and \(E\) )

\[
\begin{array}{r l} \langle x, u v \rangle = (u v) (x) & = \sum_ {i} u (y _ {i}) v (z _ {i}) = v \left(\sum_ {i} u (y _ {i}) z _ {i}\right) \\ & = v (((u \otimes 1 _ {\mathbb {E}}) \circ c) (x)) = \langle ((u \otimes 1 _ {\mathbb {E}}) \circ c) (x), v \rangle . \end{array}
\]

This can be interpreted by saying that, for all u homogeneous of degree p in \(E^{*gr}\) , the left multiplication \(e(u): v \mapsto uv\) in \(E^{*gr}\) is the graded transpose of the graded endomorphism of degree -p

\[
i (u) = (u \otimes 1 _ {\mathrm{E}}) \circ c\tag{46}
\]

of E; hence, in the above notation,

\[
(i (u)) (x) = \sum_ {i} u (y _ {i}) z _ {i}.
\]

Formula (46) also defines an element \(i(u)\in\mathrm{Endgr}_{\mathbf{A}}(\mathbf{E})\) for every element \(u\in\mathbf{E}^{*\mathbf{g}\mathbf{r}}\) ; for all \(x\in\mathbf{E}\) and all \(u\in\mathbf{E}^{*\mathbf{g}\mathbf{r}}\) , we write

\[
x \llcorner u = (i (u)) (x)\tag{47}
\]

so that, for u and v in E*gr,

\[
\langle x, u v \rangle = \langle x \llcorner u, v \rangle .
\]

The element \(x \llcorner u\) of E is called the right inner product of x by u.

Here again the element u which ``operates'' on x is placed at the free end of the horizontal line in L.

For any two elements \(u, v\) of \(\mathbf{E}^{*\mathrm{gr}}\) ,

\[
x \llcorner (u v) = (x \llcorner u) \llcorner v,\tag{48}
\]

in other words

\[
i (u v) = i (v) \circ i (u).
\]

As above let \(c(x) = \sum_{i} y_{i} \otimes z_{i}\) , so that \(x \llcorner (uv) = \sum_{i} (uv)(y_{i})z_{i}\) . If

\[
c (y _ {i}) = \sum_ {j} y _ {i j} ^ {\prime} \otimes y _ {i j} ^ {\prime \prime},
\]

then

\[
x \llcorner (u v) = \sum_ {i, j} u \left(y _ {i j} ^ {\prime}\right) v \left(y _ {i j} ^ {\prime \prime}\right) z _ {i}.\tag{49}
\]

On the other hand, if \(c(z_{i}) = \sum_{k} z_{ik}' \otimes z_{ik}''\) , then

\[
(x \llcorner u) \llcorner v = \sum_ {i, k} u (y _ {i}) v (z _ {i k} ^ {\prime}) z _ {i k} ^ {\prime \prime}.\tag{50}
\]

Now, the coassociativity of E shows that (no. 2, Proposition 1)

\[
\sum_ {i, j} y _ {i j} ^ {\prime} \otimes y _ {i j} ^ {\prime \prime} \otimes z _ {i} = \sum_ {i, k} y _ {i} \otimes z _ {i k} ^ {\prime} \otimes z _ {i k} ^ {\prime \prime}\tag{51}
\]

and the equality of expressions (49) and (50) follows from the fact that they are respectively the image of the left and right hand sides of (51) under the linear mapping f from \(E \otimes E \otimes E\) to E such that \(f(x \otimes y \otimes z) = u(x)v(y)z\) .

We recall on the other hand (no. 2, Proposition 3) that the unit element of the algebra \(\mathbf{E}^{*\mathrm{gr}}\) is the linear form \(e: x \mapsto \gamma(x).1\) ; hence

\[
x \llcorner e = \sum_ {i} \gamma (y _ {i}) z _ {i} = x
\]

by virtue of the definition of counit. As the mapping \(u \mapsto i(u)\) is linear, it is seen that the external law of composition \((u, x) \mapsto x \llcorner u\) defines a right E*gr-module structure on E.

Similarly we define, for all \(u \in \mathrm{E}^{*\mathbf{g}\mathbf{r}}\) , the endomorphism of \(\mathbf{E}\)

\[
i ^ {\prime} (u) = \left(1 _ {E} \otimes u\right) \circ c\tag{52}
\]

and, for all \(x \in \mathbf{E}\) , we write

\[
\left(i ^ {\prime} (u)\right) (x) = u \lrcorner x\tag{53}
\]

and this element of E is called the left inner product of x by u. As above it is seen that the external law \((u, v) \mapsto u \lrcorner x\) defines a left \(\mathrm{E}^{*gr}\) -module structure on E. Moreover, these two structures are compatible, in other words,

\[
(u \lrcorner x) \llcorner v = u \lrcorner (x \llcorner v)\tag{54}
\]

for \(u, v\) in E*gr (II, § 1, no. 14). With the same notation as above, the left hand side of (54) is \(\sum_{i,j} u(z_i) v(y_{ij}') y_{ij}''\) and the right hand side is \(\sum_{i,k} v(y_i) u(z_{ik}') z_{ik}''\) ; their equality follows from the fact that they are the respective images of the left and right hand sides of (51) under the linear mapping \(g\) of E ⊗ E ⊗ E into E such that \(g(x \otimes y \otimes z) = v(x) u(z)y\) .

It is therefore seen that the two external laws of composition on E defines on this set an \((\mathrm{E}^{*}\mathrm{gr}, \mathrm{E}^{*}\mathrm{gr})\) -bimodule structure.

When the cogebra E is cocommutative, then \(u \lrcorner x = x \llcorner u\) for all \(x \in \mathbf{E}\) and \(u \in \mathrm{E}^{*\mathrm{gr}}\) ; when it is anticocommutative (§ 4, no. 9) and \(u \in \mathrm{E}_p^*\) and \(x \in \mathrm{E}_q\) , we can write \(c(x) = \sum_{0 \leqslant j \leqslant q} \left( \sum_i y_{ij} \otimes z_{i,q-j} \right)\) with \(y_{ij}\) and \(z_{ij}\) in \(\mathrm{E}_j\) for all \(j\) and then by hypothesis

\[
\sum_ {i} z _ {i j} \otimes y _ {i, q - j} = (- 1) ^ {j (q - j)} \sum_ {i} y _ {i j} \otimes z _ {i, q - j}.
\]

By definition, \(x \llcorner u = \sum_{0 \leqslant j \leqslant q} \left( \sum_{i} u(y_{ij}) z_{i,q-j} \right)\) and

\[
u \lrcorner x = \sum_ {0 \leqslant j \leqslant q} \left(\sum_ {i} u (z _ {i, q - j}) y _ {i j}\right).
\]

As \(u(y_{ij}) = 0\) (resp. \(u(z_{i,q-j}) = 0\) unless \(j = p\) (resp. \(q - j = p\) ), it is seen by the above that \(u \lrcorner x = (-1)^{p(q-p)}x \llcorner u\) .

Finally, let \(\phi: \mathrm{E} \to \mathrm{F}\) be a graded cogebra morphism; then it has been seen (no. 2) that the graded transpose \(^t\phi: \mathrm{F}^{*\mathrm{gr}} \to \mathrm{E}^{*\mathrm{gr}}\) is a graded algebra homomorphism; therefore, for \(x \in \mathrm{E}, u, v\) in \(\mathrm{F}^{*\mathrm{gr}}\) ,

\[
\begin{array}{r l} \langle \phi (x \llcorner {} ^ {t} \phi (u)), v \rangle & = \langle x \llcorner {} ^ {t} \phi (u), ^ {t} \phi (v) \rangle = \langle x, ^ {t} \phi (u) ^ {t} \phi (v) \rangle = \langle x, ^ {t} \phi (u v) \rangle \\ & = \langle \phi (x), u v \rangle = \langle \phi (x) \llcorner u, v \rangle \end{array}
\]

whence

\[
\phi (x) \llcorner u = \phi (x \llcorner {} ^ {t} \phi (u));\tag{55}
\]

and similarly

\[
u \lrcorner \phi (x) = \phi (^ {t} \phi (u) \lrcorner x).\tag{56}
\]

In other words, \(\phi\) is an \(\mathbf{F}^{*\mathrm{gr}}\) -homomorphism of the \((\mathrm{E}^{*\mathrm{gr}}, \mathrm{E}^{*\mathrm{gr}})\) -bimodule E into the \((\mathrm{F}^{*\mathrm{gr}}, \mathrm{F}^{*\mathrm{gr}})\) -bimodule F, relative to the ring homomorphism

\[
{ } ^ { t } \phi : \mathrm{F} ^ { * \mathrm{gr} } \rightarrow \mathrm{E} ^ { * \mathrm{gr} } .
\]

Examples. In particular the above can be applied when E is one of the graded cogebras \(\mathsf{T}(\mathbf{M})\) , \(\mathsf{S}(\mathbf{M})\) or \(\bigwedge(\mathbf{M})\) for an A-module M (no. 1, Examples 5, 6 and 7). To find explicitly the bimodule structures thus obtained, we again identify a homogeneous element f of degree n in \(\mathsf{T}(\mathbf{M})^{*\mathbf{g}\mathbf{r}}\) (resp. \(\mathsf{S}(\mathbf{M})^{*\mathbf{g}\mathbf{r}}\) , resp. \(\bigwedge(\mathbf{M})^{*\mathbf{g}\mathbf{r}}\) ) with an n-linear form (resp. symmetric n-linear form, resp. alternating n-linear form, also called an n-form) on \(M^{n}\) . It suffices to express the products

\[
\left(x _ {1} \otimes x _ {2} \otimes \dots \otimes x _ {p}\right) \llcorner f (\text { resp. } (x _ {1} x _ {2} \dots x _ {p}) \llcorner f,
\]

resp. \((x_{1} \wedge x_{2} \wedge \cdots \wedge x_{p}) \llcorner f)\) for every finite sequence \((x_{i})_{1 \leqslant i \leqslant p}\) of elements of M and the analogues for the left inner product. Now, definitions (46) and (52) and formulae (3), (6) and (9) of no. 1 give respectively:

\[
\left\{ \begin{array}{l} (x _ {1} \otimes x _ {2} \otimes \dots \otimes x _ {p}) \llcorner f = f \lrcorner (x _ {1} \otimes x _ {2} \otimes \dots \otimes x _ {p}) = 0 \\ (x _ {1} x _ {2} \ldots x _ {p}) \llcorner f = f \lrcorner (x _ {1} x _ {2} \ldots x _ {p}) = 0 \quad \text {for} \quad p <   n. \\ (x _ {1} \wedge x _ {2} \wedge \dots \wedge x _ {p}) \llcorner f = f \lrcorner (x _ {1} \wedge x _ {2} \wedge \dots \wedge x _ {p}) = 0 \end{array} \right.\tag{57}
\]

For \(p \geqslant n\) , we have respectively

\begin{enumerate}
\def\labelenumi{(\arabic{enumi})}
\setcounter{enumi}{57}
\tightlist
\item
\end{enumerate}

\[
\left(x _ {1} \otimes x _ {2} \otimes \dots \otimes x _ {p}\right) \llcorner f = f \left(x _ {1}, \dots , x _ {n}\right) x _ {n + 1} \otimes \dots \otimes x _ {p}\tag{59}
\]

\[
\left(x _ {1} x _ {2} \dots x _ {p}\right) \llcorner f = \sum_ {\sigma} f \left(x _ {\sigma (1)}, \dots , x _ {\sigma (n)}\right) x _ {\sigma (n + 1)} \dots x _ {\sigma (p)}\tag{60}
\]

\[
\left(x _ {1} \wedge x _ {2} \wedge \dots \wedge x _ {p}\right) \llcorner f = \sum_ {\sigma} \varepsilon_ {\sigma} f \left(x _ {\sigma (1)}, \dots , x _ {\sigma (n)}\right) x _ {\sigma (n + 1)} \wedge \dots \wedge x _ {\sigma (p)}
\]

(where, in (59) and (60), the summations are taken over the permutations \(\sigma \in \mathfrak{S}_p\) which are increasing on each of the intervals \([1, n]\) and \([n + 1, p]\) of \(\mathbf{N}\) ); and similarly

\begin{enumerate}
\def\labelenumi{(\arabic{enumi})}
\setcounter{enumi}{60}
\tightlist
\item
\end{enumerate}

\[
f \lrcorner (x _ {1} \otimes x _ {2} \otimes \dots \otimes x _ {p}) = f (x _ {p - n + 1}, \dots , x _ {p}) x _ {1} \otimes x _ {2} \otimes \dots \otimes x _ {p - n}\tag{62}
\]

\[
f \lrcorner (x _ {1} x _ {2} \dots x _ {p}) = \sum_ {\sigma} f (x _ {\sigma (p - n + 1)}, \dots , x _ {\sigma (p)}) x _ {\sigma (1)} \dots x _ {\sigma (p - n)}\tag{63}
\]

\[
f \lrcorner (x _ {1} \wedge x _ {2} \wedge \dots \wedge x _ {p}) =
\]

\[
\sum_ {\sigma} \varepsilon_ {\sigma} f (x _ {\sigma (p - n + 1)}, \dots , x _ {\sigma (p)}) x _ {\sigma (1)} \wedge \dots \wedge x _ {\sigma (p - n)}
\]

(where, in (62) and (63), the summations are taken over the permutations \(\sigma \in \mathfrak{S}_p\) which are increasing on each of the intervals \([1, p - n]\) and \([p - n + 1, p]\) of \(\mathbf{N}\) ).

